\subsection{Query and shared alignment sensitivity}
The accuracy comparison uses the same automatically refined current point for every route. Of 309 screening tracks, 47 current queries have error above 3 mm. Thirty-two of these are among the 111 tracks unreachable by Original 7. Diverse search and the full spline hit 178 and 157 of 194 queries within 1.5 mm, 57 and 46 of 68 queries in the 1.5--3 mm band, and 9 and 8 of 47 above 3 mm. Holding the crop and anatomical costs fixed and replacing the query with the annotated current mark changes Original 7 from 184 to 211 and the spline from 211 to 239 hits. Diverse search changes from 244 to 243. The advantage in the automatic-query setting narrows when the exact current mark is supplied.

Some staging failures arise before local search. Among supported crops, 41 earlier reference points lie more than 3 mm beyond the entire earlier crop in the shared frame. In two patient pairs this distance exceeds 100 mm for 28 tracks. Audit of 58 lesion marks against the masks in both visits verifies the coordinate convention in these pairs. The shared affine leaves these targets outside local support. Both Diverse search and the spline return 28 errors above 200 mm. Their available-point median errors are 1.68 and 2.41 mm and their p90s are 30.40 and 29.26 mm, respectively. All tracks remain in the accuracy denominators.
